\chapter[Functional-Analytic Properties]{Functional-Analytic Properties of the Universal Optimality Defect}\label{ch:sec:functional-analytic-properties-of-the-universal-optimal}
\section{Introduction}\label{ch:sec:introduction-analytic_uod-4}
This chapter studies the global analytical properties of the Universal
Optimality Defect (UOD). Throughout,
\[
F=Q\circ L:\mathcal P^\circ\rightarrow F(\mathcal P^\circ),
\qquad \mathcal D(P,Q)=\|F(P)-F(Q)\|^2,
\]
where \(\mathcal P^\circ\) denotes the interior of the
probability simplex. Part~\ref{ch:part:posterior} proves that \(F\) is a global
diffeomorphism onto its image. The results below are consequences of this fact and of the smoothness of the logarithm on $(0,\infty)$.
\section{Smoothness}\label{ch:sec:smoothness}
\begin{theorem}[Smoothness]
\label{ch:thm:23-1}
\[
\mathcal D:\mathcal P^\circ\times\mathcal P^\circ\rightarrow\mathbb R
\]
is \(C^\infty\).
\end{theorem}
\begin{proof} The logarithm is smooth on \((0,\infty)\), \(Q\) is
linear, and the squared Euclidean norm is polynomial. Therefore
\(\mathcal D\) is a composition of smooth maps. 
\end{proof}
\section{Positivity and Separation}\label{ch:sec:positivity-and-separation}
\begin{theorem}[Positivity]
\label{ch:thm:23-2}
$\mathcal D(P,Q)\ge0,$
with equality iff \(P=Q\).
\end{theorem}
\begin{proof} Non-negativity follows from the Euclidean norm. If
\(\mathcal D(P,Q)=0\), then \(F(P)=F(Q)\). By the results of Part~\ref{ch:part:posterior}, \(F\) is
injective, hence \(P=Q\). The converse is immediate. 
\end{proof}
\section{Continuity}\label{ch:sec:continuity}
\begin{theorem}[Continuity]
\label{ch:thm:23-3}
If \(P_n\to P\) and \(Q_n\to Q\), then
$\mathcal D(P_n,Q_n)\to\mathcal D(P,Q).$
\end{theorem}
\begin{proof} By Theorem~\ref{ch:thm:23-1} the UOD is \(C^\infty\), hence in particular continuous: if \(P_n\to P\) and \(Q_n\to Q\), then \(F(P_n)\to F(P)\) and \(F(Q_n)\to F(Q)\), and the squared norm is continuous, so \(\mathcal D(P_n,Q_n)\to\mathcal D(P,Q)\).
\end{proof}
\section{Local Lipschitz Continuity}\label{ch:sec:local-lipschitz-continuity}
\begin{theorem}[Local Lipschitz]
\label{ch:thm:23-4}
For every compact set \(K\subset\mathcal P^\circ\), there exists
\(L_K>0\) such that
\[ |\mathcal D(P_1,Q_1)-\mathcal D(P_2,Q_2)|
\le L_K(\|P_1-P_2\|+\|Q_1-Q_2\|) \]
for all \(P_i,Q_i\in K\).
\end{theorem}
\begin{proof} Since \(\mathcal D\) is \(C^{1}\), its differential is
bounded on \(K\times K\). The multivariable mean-value theorem
gives the estimate. 
\end{proof}
\section{An Auxiliary Linear Lemma}\label{ch:sec:an-auxiliary-linear-lemma}
The proof of properness relies on a purely linear fact.
\begin{lemma}[Growth in the Quotient Space]
\label{ch:lem:23-5}
Let
\[ Q(a)=a-\bar a\mathbf 1, \qquad
\bar a=\frac1n\sum_i a_i. \]
Suppose a sequence \(a^{(k)}\in\mathbb R^n\) satisfies
\[ \operatorname{diam}(a^{(k)}) = \max_i
a_i^{(k)}-\min_i a_i^{(k)}
\longrightarrow\infty. \]
Then
$\|Q(a^{(k)})\|\longrightarrow\infty.$
\end{lemma}
\begin{proof} The vector \(Q(a)\) is obtained by subtracting the mean from
every coordinate. Subtracting a constant does not change pairwise
differences:
$(Q(a))_i-(Q(a))_j=a_i-a_j.$
Hence
\[ \operatorname{diam}(Q(a)) = \operatorname{diam}(a). \]
For every zero-mean vector \(x\),
\[ \|x\| \ge \frac{\operatorname{diam}(x)}{\sqrt 2}, \]
because two coordinates differ by at most \(\sqrt 2\|x\|\).
Therefore,
\[ \|Q(a^{(k)})\| \ge
\frac{\operatorname{diam}(a^{(k)})}{\sqrt 2}
\longrightarrow\infty. \]
\end{proof}
\section{Properness}\label{ch:sec:properness}
\begin{theorem}[Properness]
\label{ch:thm:23-6}
Fix \(Q\in\mathcal P^\circ\). If \(P_n\) approaches
the boundary of the simplex, then
$\mathcal D(P_n,Q)\to\infty.$
\end{theorem}
\begin{proof} Approaching the boundary implies that at least one coordinate
satisfies
$P_{n,i}\to0,$
hence
$\log\frac{P_{n,i}}{Q_i}\to-\infty.$
Since \(\sum_i P_{n,i}=1\), at least one coordinate remains bounded
away from zero along a subsequence. Therefore the diameter of the
logarithmic ratio vector
$a^{(n)}=\log P_n-\log Q$
tends to infinity. Lemma~\ref{ch:lem:23-5} yields
$\|Q(a^{(n)})\|\to\infty,$ which is precisely \(\mathcal D(P_n,Q)\to\infty\). 
\end{proof}
\section{Compact Sublevel Sets}\label{ch:sec:compact-sublevel-sets}
\begin{corollary}[Compact Sublevel Sets]
\label{ch:cor:23-7}
$S_c=\{P\in\mathcal P^\circ:\mathcal D(P,Q)\le c\}$
is compact in the relative topology of the simplex.
\end{corollary}
\begin{proof} Continuity implies that \(S_c\) is closed. Properness prevents
sequences in \(S_c\) from approaching the boundary. Since the simplex is
compact, \(S_c\) is compact. 
\end{proof}
\section{Topological Equivalence}\label{ch:sec:topological-equivalence}
\begin{theorem}[Topological Equivalence]
\label{ch:thm:23-8}
The topology induced by the UOD coincides with the manifold topology on
\(\mathcal P^\circ\).
\end{theorem}
\begin{proof} Part~\ref{ch:part:posterior} proves that \(F\) is a homeomorphism onto its image.
Since the UOD is the squared Euclidean distance in that image,
convergence in UOD is equivalent to Euclidean convergence of \(F(P)\),
hence to manifold convergence. 
\end{proof}
\section{Summary}\label{ch:sec:summary}
The Universal Optimality Defect is:
\begin{itemize}
\item infinitely differentiable;
\item positive definite;
\item jointly continuous;
\item locally Lipschitz on compact subsets;
\item proper on the interior simplex;
\item endowed with compact sublevel sets;
\item topologically equivalent to the intrinsic WIS geometry.
\end{itemize}
These properties establish the UOD as a globally well-behaved analytical
functional, providing the foundation for the comparison theorems
developed in Part~\ref{ch:part:variational}.
